\scanpage{59}\mspage{8}
% archive foliation: 58
propre et lisse sur $S$ à fibres géom. connexes de dim. $1$, genre $1$, et $\sigma$ un autom induisant $-\operatorname{id}$ sur $\underline{\operatorname{Pic}}^0_{E/S}$, \uline{équivaut} à celle \inserted{de $(E,a)$, où $a$ section de $\underline{\operatorname{Pic}}^2_{E/S}$ $(\simeq E\wedge_{E_0}E)$, ou encore à celle} \deleted{de} $(E_0,\underline I)$, où $E_0$ est une courbe elliptique relative sur $S$, et où $\underline I$ est un torseur sous \deleted{${}_2E_0\overset{\mathrm{déf}}=$} $\underline V\overset{\mathrm{déf}}={}_2(E_0)$. On aura
\[
\tag{21}\left\{\begin{array}{l}
E=\underline I\wedge_{\underline V}E_0\\
\sigma\text{ déduit de la symétrie }x\mapsto-x\text{ de }E_0\\
\qquad\text{via la formule précédente}
\end{array}\right.
\]
On aura une opération de $\underline V\unsure{\Leftarrow|=}{}_2E_0\subset E_0$ sur $E$, et un isom can.
\[
\tag{22}\begin{array}{ccc}
E/\underline V&\simeq&E_0\\
\rotatebox{90}{$\simeq$}&&\\
E\wedge_{E_0}E\simeq\underline{\operatorname{Pic}}^2_{E/S}&&
\end{array}
\]
\margin{\uline{NB} L'opération de $\underline V$ sur $E$ commute à l'action de $\sigma$}
isom qui commute aux symétries $\sigma$ sur $E$, $\sigma_0:x\mapsto-x$ sur $E_0$.

Considérons les schémas quotients
\[
\tag{23}X=E/\{1,\sigma\},\qquad\Sigma=E_0/\{1,\sigma_0\}
\]
ce sont des fibrés en \deleted{\unsure{droites}} \inserted{droites} projectives,
\deleted{\unsure{et on a les}} $V$ opère \inserted{(monomorphiquement)} sur $X$ \deleted{et sur $\Sigma$} via son \deleted{\unsure{(24)}} action sur $E$, et on aura
\[
\tag{24}X/V\simeq\Sigma,
\]
d'ailleurs le diviseur de ramification pour
\[
\tag{25}E\longrightarrow X\qquad(\text{rev. quadratique})
\]
